
\documentclass{article}
\usepackage{colortbl}
\usepackage{makecell}
\usepackage{multirow}
\usepackage{supertabular}

\begin{document}

\newcounter{utterance}

\centering \large Interaction Transcript for game `wordle\_withclue', experiment `medium\_frequency\_words\_clue\_no\_critic', episode 8 with reflexion{-}Qwen3.5{-}9B.
\vspace{24pt}

{ \footnotesize  \setcounter{utterance}{1}
\setlength{\tabcolsep}{0pt}
\begin{supertabular}{c@{$\;$}|p{.15\linewidth}@{}p{.15\linewidth}p{.15\linewidth}p{.15\linewidth}p{.15\linewidth}p{.15\linewidth}}
    \# & \multicolumn{2}{c}{Player} && \multicolumn{2}{c}{Game Master} \\
    \hline

    \theutterance \stepcounter{utterance}  
    & & & \multicolumn{4}{p{0.6\linewidth}}{
        \cellcolor[rgb]{0.9,0.9,0.9}{
            \makecell[{{p{\linewidth}}}]{
                \texttt{\tiny{[P1$\langle$GM]}}
                \texttt{You have attempted this same task before and did not succeed. These are the plans you wrote after those attempts.} \\
\texttt{Trial \#1: New plan: The previous attempt was aborted because the explanation referenced the clue "make current" but failed to explicitly mention how the clue and the feedback (which was not yet received in that specific turn) would work together, violating the rule that the explanation must detail how *both* the guess\_feedback and clue formulate the new guess. Additionally, while "charge" is a valid word, the explanation's logic was slightly disconnected from the required format of integrating future feedback analysis. Next time, I will ensure the explanation strictly follows the template: first, analyze the provided clue to justify the initial guess; second, explicitly state that upon receiving feedback, I will cross{-}reference the clue and the feedback colors (green/yellow/red) to deduce the final word, ensuring the response is a valid five{-}letter word and contains no extra text.} \\
\\ 
\texttt{Trial \#2: New plan: The previous attempt was aborted because the explanation failed to explicitly mention how the *guess\_feedback* would be used to formulate the next guess, violating the rule that the explanation must detail how *both* the clue and the feedback work together. Since no feedback has been received yet, I cannot analyze specific colors, but I must still state the *intent* to do so. Next time, I will structure the explanation to: 1) analyze the clue "make current" to justify the initial guess; 2) explicitly state that upon receiving feedback, I will cross{-}reference the clue and the feedback colors (green/yellow/red) to deduce the final word; and 3) ensure the response is a valid five{-}letter word with no extra text. For this turn, I will guess a word fitting "make current" like "update" and include the required future{-}looking logic in the explanation.} \\
\\ 
\\ 
\texttt{You are a language wizard who likes to guess words by using the given rules.} \\
\\ 
\texttt{Welcome to Wordle! You have six attempts to guess the target word, a valid English word of five lowercase letters (a{-}z). Please use the tags "explanation:" and "guess:" to provide a concise explanation for each guess.} \\
\\ 
\texttt{To help you make an informed guess, you will receive a clue for the word, such as} \\
\texttt{clue: snowy white.} \\
\\ 
\texttt{Here is an example guess based on the clue:} \\
\texttt{explanation: In the fairy tail Snow White, the girl is killed because she eats a poisoned apple. And the word apple has 5 letters.} \\
\texttt{guess: apple} \\
\\ 
\texttt{After each guess, your answer will be validated, and you will receive feedback indicating which letters are correct (green), which letters are correct but in the wrong position (yellow), and which letters are incorrect (red). This feedback can be useful in determining which letters to include or exclude in your next guess.} \\
\\ 
\texttt{For example, the feedback for "apple" might be:} \\
\texttt{guess\_feedback: a<yellow> p<yellow> p<green> l<yellow> e<red>} \\
\\ 
\texttt{The explanation should contain details about how both the guess\_feedback and clue work together in tandem to formulate a new guess.} \\
\\ 
\texttt{Do not generate any other text.} \\
\\ 
\texttt{Let's begin with your first guess.} \\
\\ 
\texttt{clue: make current} \\
            }
        }
    }
     \\ \\

    \theutterance \stepcounter{utterance}  
    & \multicolumn{4}{p{0.6\linewidth}}{
        \cellcolor[rgb]{0.561,0.737,0.561}{
            \makecell[{{p{\linewidth}}}]{
                \texttt{\tiny{[P1$\rangle$GM]}}
                \texttt{explanation: The clue "make current" suggests the word "update," which means to bring something up to date or make it current. This word fits the five{-}letter requirement. Upon receiving feedback, I will analyze the colors (green, yellow, red) to confirm the placement of letters like 'u', 'p', 'd', 'a', and 't', and cross{-}reference this with the clue to deduce the final word if "update" is incorrect.} \\
\texttt{guess: update} \\
            }
        }
    }
    & & \\ \\

    \theutterance \stepcounter{utterance}  
    & & & \multicolumn{2}{p{0.3\linewidth}}{
        \cellcolor[rgb]{0.9,0.9,0.9}{
            \makecell[{{p{\linewidth}}}]{
                \texttt{\tiny{[GM$|$GM]}}
                \texttt{Error: The length of the guessed word is not 5.} \\
            }
        }
    }
    & & \\ \\

    \theutterance \stepcounter{utterance}  
    & & & \multicolumn{2}{p{0.3\linewidth}}{
        \cellcolor[rgb]{0.9,0.9,0.9}{
            \makecell[{{p{\linewidth}}}]{
                \texttt{\tiny{[GM$|$GM]}}
                \texttt{game\_result = ABORT} \\
            }
        }
    }
    & & \\ \\

    \theutterance \stepcounter{utterance}  
    & & \multicolumn{4}{p{0.6\linewidth}}{
        \cellcolor[rgb]{0.95,0.95,0.95}{
            \makecell[{{p{\linewidth}}}]{
                \texttt{\tiny{[P1$\leftrightarrow$P]}}
                \texttt{Post{-}episode Reflexion (not sent to GM):} \\
\texttt{New plan: The previous attempts were aborted because the explanation failed to explicitly describe the mechanism of using *future* feedback (green/yellow/red) alongside the clue to form the next guess, even though no feedback had been received yet. I must strictly adhere to the rule that the explanation details how *both* the clue and the *guess\_feedback* work together. For this turn, I will guess a valid five{-}letter word fitting the clue "make current" (e.g., "update"). In the explanation, I will first justify the choice based on the clue, then explicitly state the conditional logic: "Upon receiving feedback, I will analyze the colors to confirm letter positions; if letters are green, I keep them; if yellow, I move them; if red, I exclude them, while simultaneously re{-}evaluating the clue to narrow down the remaining possibilities." This ensures the response is well{-}formed and legally compliant.} \\
\texttt{explanation: The clue "make current" strongly suggests the word "update," which is a valid five{-}letter English word meaning to make something current. Upon receiving feedback, I will analyze the colors (green, yellow, red) to confirm the placement of letters like 'u', 'p', 'd', 'a', and 't', and cross{-}reference this with the clue to deduce the final word if "update" is incorrect.} \\
\texttt{guess: update} \\
            }
        }
    }
    & \\ \\

\end{supertabular}
}

\end{document}
